% \chapter{Logistique}
% \section{Objectif du volet logistique}
% L'objectif de ce volet est de décrire comment le projet fonctionne concrètement dans le temps et dans l'espace : flux de matières, opérations, ressources, stockage, distribution et activités administratives.
% L'approche vise à passer d'une simple description du produit à une vision « chaîne de valeur / business », afin d'anticiper la faisabilité opérationnelle et d'estimer un coût de revient.

% \section{Méthode utilisée : QQOQCP (structure IBE)}
% La logistique est structurée selon la méthode QQOQCP :
% \begin{itemize}
%   \item \textbf{Quoi ?} opérations et processus (inspiration BPMN : séquencement des tâches).
%   \item \textbf{Qui ?} acteurs internes/externes, responsabilités, ressources humaines.
%   \item \textbf{Où ?} lieux, partenaires, transports amont/aval.
%   \item \textbf{Quand ?} durées et fréquences, logique de planning et gestion de commande.
%   \item \textbf{Comment ?} moyens, procédures, nomenclature, gamme opératoire, qualité, emballage, déchets.
%   \item \textbf{Pourquoi ?} calcul du coût de revient, rationalisation, création de valeur client.
% \end{itemize}


% \section{Chaîne logistique imaginée (état actuel)}
% À ce stade, la chaîne logistique est décrite en postes (macro-étapes) qui couvrent l'ensemble du flux :
% \begin{enumerate}
%   \item \textbf{Commande} des intrants et consommables : silicone, textile/étiquette, électronique, batterie, capteurs, emballages (cartons, papier bulle), cupule et consommables (cartouche encre), ainsi que la réception de commande client.
%   \item \textbf{Inventaire à la réception} : réception et mise en stock des intrants fournisseurs (approvisionnement) et des données clients (commandes, informations de contact, informations sur la prothèse, les dimensions et la postion suggérée des capteurs).
%   \item \textbf{Préparation} : préparation des matériaux + préparation capteurs/cupule
%   \item \textbf{Fabrication liner} : moulage 1, séchage 1, pose réseau capteur, moulage 2, séchage 2, contrôles qualité, mise en place cupule, impression numéro de série.
%   \item \textbf{Fabrication électronique} : préparation/fabrication de la partie électronique + contrôle qualité.
%   \item \textbf{Assemblage liner--électronique--prothèse} : connexion électronique/capteurs, adaptation prothèse si nécessaire, contrôle qualité final.
%   \item \textbf{Expédition} : emballage (carton + papier bulle) et préparation d’envoi.
%   \item \textbf{Inventaire après} : mise en stock du produit fini (liner et éléments associés), suivi de paramètres (durée de vie/outillage).
%   \item \textbf{Administratif} : gestion fournisseurs/clients (factures, suivi dossiers), et fonctions support (R\&D, RH, marketing/démarchage).
% \end{enumerate}

% \begin{table}[H]
% \centering
% \small
% \begin{tabular}{p{3.5cm} S[table-format=3.2] S[table-format=3.2]}
% \toprule
% \textbf{Poste} & {\textbf{CV matériaux}} & {\textbf{MO opérateur}}\\
% \midrule
% Commande & 298.68 & 87.75\\
% Inventaire avant & 0.00 & 87.75\\
% Préparation & 0.00 & 81.25\\
% Fabrication liner & 0.00 & 146.25\\
% Fabrication électronique & 0.00 & 13.00\\
% Assemblage liner--électronique--prothèse & 0.00 & 74.75\\
% Expédition & 0.00 & 6.50\\
% Inventaire après & 0.00 & 48.75\\
% \bottomrule
% \end{tabular}
% \caption{Coûts variables --- synthèse par poste (valeurs intermédiaires).}
% \label{tab:couts_variables_poste}
% \end{table}



% \section{Synthèse coûts (données actuelles)}
% \subsection{Coût variable unitaire (CV)}
% Le fichier de calcul indique un CV total de 298.68{\euro} par unité produite (valeur actuelle de travail) pour la composante « matériaux/consommables » au poste commande.\\
% Le modèle intègre aussi des temps opérateur par poste (préparation, fabrication, assemblage, expédition, inventaires), utilisés pour estimer la main d'œuvre variable.


% \begin{table}[H]
% \centering
% \renewcommand{\arraystretch}{1.05} % ajustement de l'espacement vertical
% \setlength{\tabcolsep}{7pt}
% \begin{tabularx}{\textwidth}{|p{3.2cm}|p{3.2cm}|p{2.6cm}|X|}
% \hline
% \textbf{Composant} & \textbf{Quantité/Liner} & \textbf{Coût Unitaire} & \textbf{Source/Commentaire} \\
% \hline
% Silicone & 360g (2x180g) & 10\texteuro/g & Matière contituant le liner \\
% \hline
% Textile + étiquette & 1 pièce & 25\texteuro/pcs & Chaussettes techniques personnalisées \\
% \hline
% Électronique & 1 pièce & 200\texteuro/pcs & Carte de contrôle STM  \\
% \hline
% % Batterie & 1 pièce & 10\texteuro/pcs & Batterie rechargeable pour alimentation du système \\
% % \hline
% % Capteurs & 15 pièces & 10\texteuro/pcs & Réseau de capteurs de pression intégrés \\
% % \hline
% Cupule & 1 pièce & 30\texteuro/pcs & Point d'encrage avec la prothèse \\
% \hline
% Carton d'emballage & 1 pièce & 10\texteuro/pcs & Emballage de protection pour expédition \\
% \hline
% Papier bulle & 2 pièces & 10\texteuro/pcs & Protection liner et prothèse séparément \\
% \hline
% Matérielle de marcage & 1 pièce & 200\texteuro/pcs & Impression code UDI (usage infime par liner) \\
% \hline
% \end{tabularx}
% \caption{Nomenclature des composants et hypothèses de coûts (version intermédiaire).}
% \label{tab:composants-couts}
% \end{table}

% \newpage

% \subsection{Coûts fixes (CF) et structure}
% Les coûts fixes listés comprennent : loyer, moules, équipements (ex. dégazage, étuve, extraction, métrologie), mobilier, services (comptable, assurances, leasing) et salaires.\\
% Le tableau de synthèse des CF met en évidence :

% \begin{longtable}{p{5.2cm} S[table-format=6.2] S[table-format=4.0] p{2.6cm}}
% \toprule
% \textbf{Nom} & {\textbf{Coût unitaire}} & {\textbf{Qté}} & \textbf{Périodicité}\\
% \midrule
% \endhead
% Loyer bâtiment & 4000.00 & 1 & mensuel\\
% Moule positif petit & 1500.00 & 11 & 60 mois\\
% Moule négatif & 1500.00 & 11 & 60 mois\\
% Moule positif grand & 1500.00 & 11 & 60 mois\\
% Machine pour écrire le code UDI & 5000.00 & 1 & 60 mois\\
% Outillage en tout genre & 12500.00 & 1 & 60 mois\\
% Balance & 420.00 & 1 & 60 mois\\
% Chambre vide (dégazage) & 1000.00 & 1 & 60 mois\\
% Étuve & 1890.00 & 1 & 60 mois\\
% Duromètre & 121.00 & 3 & 60 mois\\
% Extracteur & 3000.00 & 1 & 60 mois\\
% Mobilier + accessoires & 4900.00 & 1 & 60 mois\\
% Salaire moyen ouvrier & 9056.85 & 2 & mensuel\\
% Salaire gestionnaire/commercial & 10703.55 & 1 & mensuel\\
% Comptable & 3600.00 & 1 & annuel (12 mois)\\
% Leasing & 800.00 & 1 & mensuel\\
% Assurance incendie & 300.00 & 1 & mensuel\\
% Assurance exploitation & 300.00 & 1 & mensuel\\
% \bottomrule
% \caption{Extrait de coûts fixes (valeurs de travail, à valider et documenter).}
% \end{longtable}


% \begin{itemize}
%   \item un \textbf{coût de démarrage} (investissement / mois initial) : \SI{116390.25}{\euro},
%   \item un \textbf{niveau mensuel récurrent} : \SI{34217.25}{\euro} (avec une variation à \SI{37817.25}{\euro} sur certains mois dans la table).
% \end{itemize}

% \section{Points ouverts et prochaines validations}
% Plusieurs données sont encore non acquises sur plusieurs lignes (prix fournisseurs, quantités, pertes, temps exacts), ce qui implique une phase de consolidation.
% Les priorités de validation logistique sont :
% \begin{itemize}
%   \item \textbf{Approvisionnement} : délais.
%   \item \textbf{Gamme opératoire} : stabiliser les temps (opérateur/machine), identifier les goulots (séchage, QC, assemblage).
%   \item \textbf{Qualité \& traçabilité} : formaliser contrôles, critères d’acceptation, documentation.
%   \item \textbf{Distribution \& SAV} : définir le circuit (livraison, retours, support), souvent sous-estimé dans les rapports.
%   \item \textbf{Cohérence finance--capacité} : relier volumes, capacité atelier, coûts fixes et coût unitaire.
% \end{itemize}


\chapter{Logistique}


% % Couleurs
% \definecolor{tecfitblue}{RGB}{0,70,130}
% \definecolor{lightgray}{RGB}{245,245,245}
% \definecolor{midgray}{RGB}{180,180,180}

% % Titres
% \titleformat{\section}{\large\bfseries\color{tecfitblue}}{}{0em}{\thesection.\quad}[\vspace{-0.3em}\rule{\linewidth}{0.5pt}]
% \titleformat{\subsection}{\normalsize\bfseries\color{tecfitblue}}{}{0em}{\thesubsection\quad}
% \titlespacing*{\section}{0pt}{1.5em}{0.8em}
% \titlespacing*{\subsection}{0pt}{1em}{0.4em}

% % En-tête et pied de page
% \pagestyle{fancy}
% \fancyhf{}
% \fancyhead[L]{\small\textcolor{tecfitblue}{\textbf{Rapport Logistique --- Liner Prothétique Adaptatif}}}
% \fancyhead[R]{\small\textcolor{midgray}{Projet Tecfit --- IBE 2025-2026}}
% \fancyfoot[C]{\small\thepage}
% \renewcommand{\headrulewidth}{0.4pt}
% \renewcommand{\footrulewidth}{0.4pt}

% % Hyperlinks
% \hypersetup{
%   colorlinks=true,
%   linkcolor=tecfitblue,
%   urlcolor=tecfitblue,
%   pdftitle={Rapport Logistique Tecfit},
%   pdfauthor={Projet IBE}
% }



%% ===== PAGE DE TITRE =====
% \begin{titlepage}
%   \centering
%   \vspace*{2cm}
%   {\Huge\bfseries\color{tecfitblue} Rapport de Logistique\par}
%   \vspace{0.5cm}
%   {\Large\color{midgray}\rule{10cm}{1pt}\par}
%   \vspace{0.5cm}
%   {\large Liner Prothétique Adaptatif\par}
%   \vspace{0.3cm}
%   {\normalsize Nouvelle entreprise issue du projet universitaire Tecfit\par}
%   \vspace{2cm}
%   \begin{tabular}{rl}
%     \textbf{Projet :} & Transdisciplinaire IBE 2025--2026 \\[4pt]
%     \textbf{Partenaire industriel :} & Tecfit SRL \\[4pt]
%     \textbf{Cadre financier :} & Cas D --- Nouvelle entreprise \\[4pt]
%     \textbf{Date :} & Avril 2026 \\[4pt]
%     \textbf{Méthode logistique :} & QQOQCP + BPMN \\
%   \end{tabular}
%   \vspace{3cm}
%   \begin{center}
%     \textcolor{lightgray}{\rule{14cm}{0.3pt}}\\[0.3cm]
%     {\small\textit{Ce rapport détaille la conception complète de la chaîne logistique\\
%     pour la fabrication et la distribution de liners prothétiques intelligents.}}
%   \end{center}
%   \vfill
% \end{titlepage}


%% ===== SECTION 1 =====
\section{Introduction et Contexte}

La partie logistique du projet a été conçue selon la méthode \textbf{QQOQCP}
(Quoi, Qui, Où, Quand, Comment, Pourquoi), telle que prescrite dans le powerpoint de la partie logistique (source).
Cette méthode permet de modéliser la chaîne opérationnelle complète, depuis l'approvisionnement en matières
premières jusqu'à la livraison du produit fini chez le client, en passant par toutes les étapes de
fabrication et de contrôle qualité. Le résultat est une chaîne logistique représentée sous forme de diagrammes BPMN (Business Process Model and Notation) et d'une
nomenclature (Bill of Materials).

%% ===== SECTION 2 =====
\section{Quoi : Le Produit et les Processus (BPMN)}

\subsection{Description du produit fini}

Le produit livré au client est un ensemble (P-000) comprenat :
\begin{itemize}
  \item Un liner en silicone avec réseau de capteurs intégré et cupule (référence SA1);
  \item Une partie électronique (PCB, STM32, batterie LiPo) assemblée à la prothèse;
  \item Un emballage pour l'expédition et la livraison (carton, papier bulle, étiquettes).
\end{itemize}

Le coût de revient total par liner estimé est de 354,24\euro{}, réparti en trois sous-ensembles~:

\begin{center}
\begin{tabular}{llr}
  \toprule
  \textbf{Sous-ensemble} & \textbf{Code} & \textbf{Coût/liner} \\
  \midrule
  Liner avec capteur     & SA1           & 183,82~\euro{} \\
  Électronique           & SA2           & 122,09~\euro{}       \\
  Emballage + Livraison  & SA3           & 48,32~\euro{}        \\
  \midrule
  \textbf{Total produit fini} & \textbf{P-000} & \textbf{354,24~\euro{}} \\
  \bottomrule
\end{tabular}
\end{center}
\vspace{0.5cm}

\subsection{Macro-processus opérationnel}

Le BPMN conçu pour ce projet articule cinq grands blocs fonctionnels, chacun décomposé en
sous-étapes (processus):

\begin{enumerate}
  \item Commande : déclenché mensuellement par le Gestionnaire~;
  \item Inventaire \& Stockage : contrôle qualité à l'entrée pour chaque matière/donnée;
  \item Fabrication du liner : production du liner (moulage 1 et 2, mise en pcle capteur, cupule, etc);
  \item Fabrication électronique : assemblage de la partie électronique (PCB, STM32, batterie) en parallèle du liner;
  \item Assemblage : connexion liner-électronique-prothèse et calibrage;
  \item Emballage \& Expédition : conditionnement et remise au livreur UPS.
\end{enumerate}

Chaque processus est modélisé sous forme de logigramme BPMN, avec des points de contrôle qualité formels à chaque étape critique (ex. après moulage 1, après moulage 2, après assemblage électronique, etc),
garantissant la traçabilité et la conformité réglementaire du dispositif médical de classe IIa.

%% ===== SECTION 3 =====
\section{Qui : Acteurs et Ressources Humaines}

\subsection{Rôles internes}

Deux profils principaux structurent l'organisation de l'atelier.

\textbf{Le Gestionnaire} est responsable de l'ensemble des activités de pilotage:
\begin{itemize}
  \item Gestion des commandes fournisseurs (mensuel);
  \item Réception et validation des commandes clients;
  \item Gestion RH, marketing, R\&D et administrative;
  \item Gestion des litiges fournisseurs et clients;
  \item Émission des factures.
\end{itemize}

\textbf{L'Ouvrier qualifié} exécute toutes les opérations de production et logistique physique~:
\begin{itemize}
  \item Inventaire avant production (mensuel);
  \item Préparation des matières (pesée silicone, découpe textile, vérification composants);
  \item Fabrication du liner (moulage 1 et 2, pose réseau capteur, cupule, impression UDI);
  \item Assemblage électronique et connexion au liner~;
  \item Contrôles qualité (4 points de contrôle formels dans le flux);
  \item Emballage, inventaire après production, préparation expédition.
\end{itemize}

\subsection{Dimensionnement RH}

Sur base du volume de production estimé et des temps opératoires de la gamme, l'entreprise prévoit:

\vspace{0.5em}
\begin{center}
\begin{tabular}{lrr}
  \toprule
  \textbf{Profil} & \textbf{Salaire brut mensuel} & \textbf{Coût horaire (charges incluses)} \\
  \midrule
  Gestionnaire/commercial & 5\,800\euro{} & 39,56\euro{}/h \\
  Ouvrier qualifié  & 4\,445\euro{} & 30,32~\euro{}/h \\
  \bottomrule
\end{tabular}
\end{center}
\vspace{0.5em}

Ces montants intègrent les charges patronales (25\,\%), le 13\textsuperscript{ème} mois et le pécule
de vacances, conformément à la législation sociale belge.

%% ===== SECTION 4 =====
\section{Où : Localisation, Fournisseurs et Transport}

\subsection{Site de production}

L'entreprise prévoit de louer un bâtiment industriel de 300m² en région liégeoise
(Belgique), pour un loyer mensuel estimé à \textbf{3\,500~\euro{}/mois}. Cette superficie est
dimensionnée pour accueillir l'atelier de moulage, la zone d'assemblage électronique, le stock de
matières premières et le bureau du gestionnaire.

\subsection{Réseau de fournisseurs}

% La nomenclature (BOM) identifie les fournisseurs pour chaque composant~:

% \vspace{0.5em}
% \begin{center}
% \begin{tabular}{llrr}
%   \toprule
%   \textbf{Composant} & \textbf{Fournisseur} & \textbf{Prix unitaire} & \textbf{Unité commande} \\
%   \midrule
%   Silicone             & Devis     & 52,03~\euro{}/kg & Par 10~kg   \\
%   Textile nylon-spandex & Etsy            & 12,50~\euro{}    & Par 50~pcs  \\
%   Capteur              & (déduit BOM)     & 74,89~\euro{}    & Unité       \\
%   Cupule               & PCBway           & 18,88~\euro{}    & Par 10~pcs  \\
%   PCB                  & PCBally          & 80,00~\euro{}    & Par 100~pcs \\
%   STM32 (µcontrôleur)  & Mouser.fr        & 17,11~\euro{}    & Unité       \\
%   Batterie LiPo        & Amazon.be        & 26,00~\euro{}    & Unité       \\
%   Carton d'expédition  & Rajapack.be      & 1,68~\euro{}     & Par 1\,000  \\
%   Papier bulle         & Brico.be         & 0,58~\euro{}/m²  & Par 50~m²   \\
%   Étiquettes           & Brother (QL-600) & 0,24~\euro{}     & Par 100pcs \\
%   \bottomrule
% \end{tabular}
% \end{center}
% \vspace{0.5em}

La nomenclature identifie dix composants clés avec leurs prix unitaires : silicone à 52,03\euro{}/kg, textile nylon-spandex à 12,5\euro{}/textile, capteurs à 74,89\euro{}/capteur, cupules à 18.88\euro{}/cupule, PCB à 80\euro{}/PCB, STM32 à 17,11\euro{}/STM, batterie LiPo à 26,00\euro{}/Batterie, carton d'expédition à 1,68\euro{}/carton, papier bulle à 0,58\euro{}/m², et étiquettes à 0,24\euro{}/étiquette. Le détail complet est disponible dans l'Excel en annexe.


Les commandes sont mensuelles, basées sur un calcul des quantités nécessaires pour le
mois suivant, avec un inventaire de début de mois réalisé par le gestionnaire.

\subsection{Transport et livraison}

La livraison des produits finis aux clients est faite par UPS, avec un coût de livraison de
45,34~\euro{} par colis. Ce tarif représente 12,8\,\% du coût de revient
total, ce qui en fait le poste de coût important. Le livreur
UPS passe directement à l'atelier pour prendre en charge les colis après leur mise en stock et mise à
jour d'inventaire.

%% ===== SECTION 5 =====
\section{Quand : Temporalité et Fréquences}

\subsection{Périodicité des processus}

Le flux opérationnel distingue deux fréquences fondamentales~:

\begin{itemize}
  \item Processus mensuels (pris en charge par le Gestionnaire) : commandes fournisseurs, inventaire
    général, gestion RH, marketing, R\&D, suivi comptable~;
  \item Processus par liner (à chaque commande client): préparation matières,
    fabrication, assemblage, contrôle qualité, emballage, expédition.
\end{itemize}

\subsection{Temps de cycle de fabrication}

Les temps opératoires clés identifiés dans la gamme sont:

\vspace{0.5em}
\begin{center}
\begin{tabular}{lr}
  \toprule
  \textbf{Opération} & \textbf{Durée estimé pour un opérateur} \\
  \midrule
  Préparation des matières premières     & 2h \\
  Fabrication liner                      & 2h05 ($+$ $2 \times 4$~h de polymérisation)  \\
  Assemblage électronique hors liner     & 2h     \\
  Connexion électronique + calibrage     & 1h       \\
  Emballage + inventaire + expédition    & 1/3h \\
  \bottomrule
\end{tabular}
\end{center}
\vspace{0.5em}
Le temps de cycle total par liner est dominé par les deux phases de polymérisation du silicone.
La production en batch sur les \textbf{11 jeux de moules} disponibles (positifs
petits, grands et négatifs) permet de paralléliser la fabrication.




%% ===== SECTION 6 =====
\section{Comment : Gamme Opératoire et Nomenclature}

\subsection{Nomenclature (BOM) structurée}

La nomenclature est organisée sur trois niveaux~:
\begin{itemize}
  \item Niveau 0: Produit fini P-000 (Liner + électronique emballé);
  \item Niveau 1: Trois sous-ensembles (SA1 Liner, SA2 Électronique, SA3 Emballage/Livraison)~;
  \item Niveaux 2: Composants élémentaires avec codes, prix unitaires, quantités par liner
    et sources fournisseurs.
\end{itemize}

\subsection{Gamme opératoire : Fabrication Liner}

La fabrication du liner suit une séquence stricte en deux cycles de moulage.

\textbf{Préparation:}
\begin{enumerate}
  \item Pesée du silicone ($\sim$1,5kg), mélange (3~min), dégazage sous vide (10min)~;
  \item Découpe du textile aux bonnes dimensions;
  \item Vérification et formation du réseau de capteurs selon données patient.
  \item Préparation de la cupule (choix selon prothèse, vérification dimensions).
\end{enumerate}

\textbf{Moulage 1 : Intégration textile:}
\begin{enumerate}
  \item Choix des moules (négatif + positif grand) selon morphologie du patient;
  \item Positionnement du textile au fond du moule négatif;
  \item Versage du silicone préparé dans le moule négatif;
  \item Positionnement et vissage du moule positif grand~
  \item Polymérisation : attente 4~heures;
  \item Dévissage, retrait du moule positif, contrôle qualité 1 (aspect, homogénéité,
    liaison textile-silicone).
\end{enumerate}

\textbf{Moulage 2 : Intégration réseau capteur~:}
\begin{enumerate}[leftmargin=1.5em]
  \item Positionnement du réseau de capteurs dans le moulage 1;
  \item Versage d'une nouvelle couche de silicone, pose du moule positif petit;
  \item Polymérisation: attente 4~heures;
  \item Contrôle qualité 2 (réseau capteur + liner).
\end{enumerate}

\textbf{Finitions:}
\begin{enumerate}[leftmargin=1.5em]
  \item Mise en place de la cupule, contrôle qualité 3;
  \item Impression du code UDI, marquage CE et logo (traçabilité réglementaire classe IIa).
\end{enumerate}

\subsection{Gamme opératoire : Fabrication et Assemblage Électronique}

La partie électronique (PCB + STM32 + batterie) est assemblée en parallèle de la fabrication
du liner, optimisant ainsi le temps de cycle global~:

\begin{enumerate}
  \item Vérification de tous les composants électroniques et de la charge batterie;
  \item Assemblage hors-liner de la partie électronique;
  \item Contrôle qualité électronique (fonctionnement et connectique);
  \item Modification/perçage de la prothèse pour l'intégration mécanique;
  \item Connexion du réseau de capteurs du liner à la partie électronique~;
  \item Calibrage du liner (adaptation au patient);
  \item Contrôle qualité final de l'ensemble liner-électronique-prothèse.
\end{enumerate}

En cas de non-conformité à chaque point de contrôle, un arbre de décision est prévu: modification si
possible, mise au rebut avec récupération des composants réutilisables (cupule, textile) si possible.

%% ===== SECTION 7 =====
\section{Infrastructure et Équipements}

\subsection{Équipements de production}

Les investissements principaux pour la mise en production sont les suivants:

\vspace{0.5em}
\begin{center}
\begin{tabular}{lrrr}
  \toprule
  \textbf{Équipement} & \textbf{Coût unitaire} & \textbf{Qté} & \textbf{Amort.} \\
  \midrule
  Moules positifs petits (aluminium)   & 1\,980~\euro{}  & 11 & 60~mois \\
  Moules négatifs (aluminium)          & 2\,500~\euro{}  & 11 & 60~mois \\
  Moules positifs grands (aluminium)   & 2\,500~\euro{}  & 11 & 60~mois \\
  Chambre à vide (dégazage)            & 1\,000~\euro{}  &  1 & 60~mois \\
  Étuve (finition silicone)            & 3\,865~\euro{}  &  1 & 60~mois \\
  Balance de précision                 &   420~\euro{}   &  1 & 60~mois \\
  Duromètres (dureté silicone Shore A) &   300~\euro{}   &  3 & 60~mois \\
  Extracteur de fumées                 & 1\,960~\euro{}  &  1 & 60~mois \\
  Machine impression UDI               & 5\,000~\euro{}  &  1 & 60~mois \\
  Outillage atelier complet            & 23\,607~\euro{} &  1 & 60~mois \\
  Equipement entreprise                & 14\,000~\euro{}  &  1 & 60~mois \\
   \midrule
  \bottomrule
\end{tabular}
\end{center}
\vspace{0.5em}

\subsection{Infrastructure numérique}

La chaîne logistique est supportée par un ERP/système de gestion de stock
(Microsoft Dynamics 365) permettant la traçabilité des stocks, la gestion des commandes
clients et fournisseurs, ainsi que la comptabilité. SolidWorks est utilisé pour la conception 3D des
moules et la vérification dimensionnelle des liners. Le coût annuel des licences logicielles est estimé
à 15\,264,95~\euro{}/an.

%% ===== SECTION 8 =====
\section{Pourquoi : Valeur Générée et Coût de Revient}

\subsection{Justification économique de la chaîne logistique}

La logistique conçue devrait avoir un coût de revient de 354,24~\euro{}/liner,
permettant de dégager une marge au vu du remboursement INAMI pour les liners sur mesure (2696,83~\euro{}), \url{https://www.inami.fgov.be/SiteCollectionDocuments/nomenclatureart28_1_20250701_01.pdf} (p7).
La valeur créée pour le patient va au-delà du coût matière: c'est l'adaptation
dynamique du liner, la précision du réseau de capteurs et la personnalisation morphologique qui
constituent la proposition de valeur.

\subsection{Conformité réglementaire comme élément logistique}

L'impression du code UDI (Unique Device Identifier, https://www.imdrf.org/sites/default/files/docs/imdrf/final/technical/imdrf-tech-131209-udi-guidance-140901.pdf) et du marquage CE classe IIa (https://eur-lex.europa.eu/legal-content/FR/TXT/?uri=celex:32017R0745R(01))
(selon le règlement UE 2017/745 sur les dispositifs médicaux) est intégrée directement dans la gamme
de fabrication, au même titre qu'une opération de production.

% \subsection{Gestion des non-conformités et taux de rebut}

% Le flux BPMN prévoit systématiquement des boucles de récupération à chaque point de contrôle qualité.
% La politique est la suivante~:

% \begin{itemize}[leftmargin=1.5em]
%   \item \textbf{Silicone/liner non conforme}~: mise au rebut du silicone, récupération de la cupule et
%     du textile si possible~;
%   \item \textbf{Électronique défaillante}~: identification et remplacement du composant défectueux~;
%   \item \textbf{Prothèse non conforme après assemblage}~: séparation des sous-ensembles, diagnostic de
%     la cause racine.
% \end{itemize}

% Cette approche limite les pertes financières tout en garantissant qu'aucun produit non conforme n'est
% expédié au client, ce qui est une exigence absolue pour un dispositif médical de classe IIa.

%% ===== SECTION 9 =====
\section{Conclusion}

La chaîne logistique du projet liner prothétique a été conçue selon
le cadre \textbf{QQOQCP}. Le BPMN couvre le flux opérationnel, de la commande
fournisseur à la livraison client. La gamme opératoire identifie le moulage en
silicone (2$\times$4h de polymérisation) comme le goulot d'étranglement principal, que la
production en parallèle sur \textbf{11~jeux de moules} permet de mitiger. Enfin, l'intégration de la
traçabilité réglementaire (UDI, CE IIa) directement dans le flux de production positionne l'entreprise
pour répondre aux exigences du marché des dispositifs médicaux dès le lancement, et facilite à terme
le rachat de l'entreprise par Tecfit SRL.






